% ===================================================================
%  Dodatek T4: Topologia solitonu, odbicia od ghost wall
%              i naturalne wykluczenie 4. generacji
%  TGP v1 — Sesja v41 (2026-04-03)
%  Wyniki numeryczne: ex127 (14/15 PASS) + ex128 (16/16 PASS)
%                     ex130 (11/11 PASS) + ex132 (14/14 PASS)
% ===================================================================
\section{Topologia solitonu TGP: odbicia od \emph{ghost wall}
  i~naturalne wykluczenie czwartej genera\-cji}%
\label{app:T4-bounce}

Niniejszy dodatek bada topologiczn\k{a} struktur\k{e} soliton\'ow TGP
na drabinie~$\varphi$ oraz wykazuje, \.ze mechanizm masy $m \propto
A_{\mathrm{tail}}^4$ \textbf{wyklucza czwart\k{a} genera\-cj\k{e}}
wewn\k{e}trznie --- niezale\.znie od ogranicze\'n LEP.
Wyniki zosta\l{}y zweryfikowane w~\texttt{ex127} (14/15~PASS).

% -------------------------------------------------------------------
\subsection{Mechanizm odbicia od \emph{ghost wall}}
\label{app:T4-bounce-mechanism}
% -------------------------------------------------------------------

Profil solitonu $g(r)$ spe\l{}nia ODE (sek.~\ref{sec:soliton}):
\[
  f(g)\,g'' + \frac{\alpha}{g}(g')^2 - f(g)\frac{2}{r}g' = V'(g),
  \qquad
  f(g) = 1 + 2\alpha\ln g,
\]
z~warunkiem pocz\k{a}tkowym $g(0) = g_0$, $g'(0)=0$.
Dla $g$ malejacego do warto\'sci \emph{ghost}: $g_* = e^{-1/(2\alpha)}$
mamy $f(g_*) = 0$, co powoduje \emph{elastyczne odbicie} profilu
(zmian\k{e} znaku $g'$).

\begin{definition}[Liczba odbi\'{c} $n_{\mathrm{bounce}}(g_0)$]%
\label{def:T4-nbounce}
$n_{\mathrm{bounce}}(g_0)$ to liczba krotno\'sci,
z~jak\k{a} $g(r)$ osi\k{a}ga $g_*$ dla $r \in [0, R_{\max}]$
przed ustabilizowaniem si\k{e} asymptotyki~$g(r) \to 1$.
\end{definition}

% -------------------------------------------------------------------
\subsection{Bounce-Generation Hypothesis (BGH)}
\label{app:T4-bgh}
% -------------------------------------------------------------------

\begin{hypothesis}[BGH --- wersja monotoniczna]\label{hyp:T4-bgh}
Wzdlu\.z drabiny $\varphi$:
$g_0^{(k)} = \varphi^{k-1} g_0^e$, $k = 1, 2, 3, 4$,
liczba odbi\'{c} $n_{\mathrm{bounce}}(g_0^{(k)})$ jest
\emph{\'scis\l{}e rosn\k{a}ca}:
\[
  n_{\mathrm{bounce}}(g_0^{(1)})
  < n_{\mathrm{bounce}}(g_0^{(2)})
  < n_{\mathrm{bounce}}(g_0^{(3)})
  < n_{\mathrm{bounce}}(g_0^{(4)}).
\]
\end{hypothesis}

\begin{remark}[Weryfikacja numeryczna (\texttt{ex127})]
Dane:
\begin{center}
\renewcommand{\arraystretch}{1.2}
\begin{tabular}{@{}ccccccc@{}}
\toprule
$k$ & Lepton & $g_0$ & $n_{\mathrm{bounce}}$ & $A_{\mathrm{tail}}$
  & RMSE$/A$ & Status \\
\midrule
1 & $e$   & 1{,}249 & 0 & 0{,}299 & 0{,}5\%  & \statuslabel{Stabilny} \\
2 & $\mu$ & 2{,}021 & 1 & 1{,}133 & 2{,}2\%  & \statuslabel{Stabilny} \\
3 & $\tau$& 3{,}189 & 3 & 2{,}295 & 8{,}0\%  & \statuslabel{Graniczny} \\
4 & $L_4$ & 5{,}291 & 6 & 3{,}771 & 36{,}4\% & \statuslabel{Zdegradowany} \\
\bottomrule
\end{tabular}
\end{center}
BGH (wersja monotoniczna) \textbf{potwierdzona} (\texttt{ex127}, B9~PASS).
Prosta wersja $n_{\mathrm{bounce}} = k-1$ nie zachodzi: $\tau$ ma 3, nie~2 odbicia.
\end{remark}

% -------------------------------------------------------------------
\subsection{RMSE-Cutoff Hypothesis (RCH): naturalne wykluczenie $k=4$}
\label{app:T4-rch}
% -------------------------------------------------------------------

Jako\'s\'c dopasowania ogona asymptotycznego mierzono przez
$\mathrm{RMSE}/A_{\mathrm{tail}}$ (wzgl\k{e}dny b\l{}\k{a}d
\'srednio-kwadratowy dopasowania $(g-1)\cdot r \approx B\cos r + C\sin r$).

\begin{hypothesis}[RCH --- kryterium jako\'sci ogona]\label{hyp:T4-rch}
Mechanizm masy $m_k \propto A_{\mathrm{tail}}(g_0^{(k)})^4$ jest
\emph{wiarygodny} tylko wtedy, gdy RMSE$/A_{\mathrm{tail}} < 10\%$.
Dla $k = 4$ ($L_4$): RMSE$/A = 36{,}4\% > 15\%$ --- ogon
asymptotyczny jest zdegradowany przez wielokrotne odbicia.
\end{hypothesis}

\begin{theorem}[RCH wyklucza czwart\k{a} genera\-cj\k{e}]%
\label{thm:T4-rch}
Niech $g_0^{(\max)}(\varepsilon)$ oznacza maksymalne $g_0$ takie, \.ze
RMSE$/A_{\mathrm{tail}} < \varepsilon$.  Dla $\varepsilon = 10\%$:
\[
  g_0^{(3)} = 3{,}189
  \;<\; g_0^{(\max)}(10\%) = 3{,}591
  \;<\; g_0^{(4)} = 5{,}291.
\]
W konsekwencji: $\tau$ jest \emph{ostatni\k{a} genera\-cj\k{a}} z~dobrze
okre\'slonym mechanizmem masowym~$A_{\mathrm{tail}}^4$.
\end{theorem}

\begin{proof}
Bezpo\'srednia weryfikacja numeryczna, \texttt{ex127}, test~B8 i~B14.
Skan g\k{e}sty $g_0 \in [1{,}05, 6{,}00]$ (50 punkt\'ow) potwierdza, \.ze
RMSE$/A < 10\%$ dla $g_0 \leq 3{,}59$ i~RMSE$/A > 10\%$ dla
$g_0 > 3{,}59$.
\end{proof}

\begin{corollary}[T-OP3 cz\k{e}\'sciowo zamkni\k{e}ty]\label{cor:T4-top3}
TGP wyklucza czwart\k{a} genera\-cj\k{e} leptonow \emph{wewn\k{e}trznie}:
bez konieczno\'sci odwo\l{}ania do danych LEP, lecz przez kryterium
stabilno\'sci asymptotyki solitonu.  W po\l{}\k{a}czeniu z~wynikami
ex126 (T-OP1~$\equiv$~T-OP3):
\begin{equation}\label{eq:T4-chain}
  \underbrace{N_{\mathrm{gen}} = 3}_{\text{RCH (T4)}}
  \;\Longrightarrow\;
  \underbrace{r = \sqrt{N-1} = \sqrt{2}}_{\text{ex126}}
  \;\Longrightarrow\;
  \underbrace{Q_K = \frac{2N}{N+1} = \frac{3}{2}}_{\text{ex118--ex119}}.
\end{equation}
\end{corollary}

% -------------------------------------------------------------------
\subsection{Mapa skanowania $g_0$: progi i~dips}
\label{app:T4-scan}
% -------------------------------------------------------------------

Skan $g_0 \in [1{,}05, 6{,}00]$ (50 punkt\'ow, \texttt{ex127}) ujawnia:

\begin{enumerate}
  \item \textbf{Monotoniczno\'s\'c $n_{\mathrm{bounce}}(g_0)$}:
    $n_{\mathrm{bounce}}$ jest niemalejace --- B5 PASS, 0/49 krok\'ow
    niegodnych.

  \item \textbf{Progi $g_0^*$} (przej\'scia $n \to n{+}1$):
    $g_0^* \approx 1{,}625$ ($0\to1$),
    $g_0^* \approx 2{,}232$ ($1\to2$),
    $g_0^* \approx 2{,}875$ ($2\to3$).

  \item \textbf{Lokalne minima $A_{\mathrm{tail}}$ przy progach}:
    W~s\k{a}siedztwie ka\.zdego progu $n_{\mathrm{bounce}}$
    wyst\k{e}puje lokalne minimum $A_{\mathrm{tail}}$ z~g\l{}\k{e}boko\'sci\k{a}
    $10{-}20\%$ (B10~PASS).  Mo\.ze to wskazywa\'c na quasi-degeneracj\k{e}
    faz soliton\'ow o~r\'o\.znej liczbie odbi\'c.

  \item \textbf{Ci\k{a}g\l{}o\'s\'c $A_{\mathrm{tail}}$}:
    Skok $A_{\mathrm{tail}}$ przy progu $n=0\to1$: $4{,}5\%$ (B11~PASS) ---
    brak nieciag\l{}o\'sci.
\end{enumerate}

% -------------------------------------------------------------------
\subsection{Degradacja RMSE i~nowe pytanie T-OP5}
\label{app:T4-top5}
% -------------------------------------------------------------------

\begin{observation}[Progresja RMSE]\label{obs:T4-rmse-prog}
Wzgl\k{e}dny b\l{}\k{a}d RMSE$/A$ gwa\l{}townie ro\'snie wzdlu\.z
drabiny~$\varphi$:
\begin{center}
\begin{tabular}{ccccc}
\toprule
Gen. & $e$ & $\mu$ & $\tau$ & $L_4$ \\
\midrule
RMSE$/A$ & 0{,}5\% & 2{,}2\% & 8{,}0\% & 36{,}4\% \\
\bottomrule
\end{tabular}
\end{center}
Wzrost w~przybli\.zeniu geometryczny z~czynnikiem ${\sim}3{-}4$ na krok.
\end{observation}

\begin{remark}[Spowalnianie stosunk\'ow $A_{\mathrm{tail}}$]
Mimo rosn\k{a}cego RMSE, warto\'sc $A_{\mathrm{tail}}$ r\'ownie\.z ro\'snie,
lecz z~\emph{malej\k{a}cym} wsp\'o\l{}czynnikiem wzrostu:
\[
  \frac{A_\mu}{A_e} = 3{,}79,\quad
  \frac{A_\tau}{A_\mu} = 2{,}03,\quad
  \frac{A_{L_4}}{A_\tau} = 1{,}64.
\]
Ta \emph{saturacja} amplitudy wraz z~degradacj\k{a} RMSE stanowi
dwie niezale\.zne przes\l{}anki wykluczenia $k = 4$ przez TGP.
\end{remark}

\begin{proposition}[T-OP5 --- cz\k{e}\'sciowe zamkni\k{e}cie, \texttt{ex130}]%
\label{prop:T4-top5}
Gęsty skan $g_0 \in [1{,}05,\,6{,}00]$ (100~punkt\'ow, \texttt{ex130})
ujawnia regularn\k{a} siatk\k{e} prog\'ow $g_0^*(k\to k{+}1)$:
\begin{equation}\label{eq:T4-bounce-thresholds}
  g_0^*(k) \;\approx\; 0{,}675\,k + 1{,}582,
  \quad k = 0,1,2,\ldots
\end{equation}
z~RMSE dopasowania $= 0{,}046$.
Przy ka\.zdym progu RMSE$/A$ ro\'snie (7/7~przej\'s\'c, E6~PASS).

Warto\'sci progowe:
\[
  g_0^*(0{\to}1) = 1{,}650,\;
  g_0^*(1{\to}2) = 2{,}250,\;
  g_0^*(2{\to}3) = 2{,}900,\;
  g_0^*(3{\to}4) = 3{,}549.
\]
Ponadto $|g_0^{(\max)}(10\%) - g_0^*(3{\to}4)| = 0{,}147 < 0{,}15$
(E3~PASS): \textbf{$g_0^{(\max)}$ le\.zy tu\.z powy\.zej progu
przej\'scia $n_{\mathrm{bounce}}$ $3{\to}4$}.
\end{proposition}

\begin{proposition}[T-OP5 --- zamkni\k{e}cie analityczne, \texttt{ex132}]%
\label{prop:T4-top5-analytic}
Nachylenie $\Delta g_0^*$ jest wyznaczone przez trzy elementy ODE solitonu TGP:
\begin{enumerate}
  \item \textbf{Cz\k{e}sto\'s\'c liniowych oscylacji} przy $g \approx 1$:
    Linearyzacja~ODE ($h = g-1$, $r\to\infty$):
    \begin{equation}\label{eq:T4-linearized}
      h'' - \frac{2}{r}\,h' + h = 0,
    \end{equation}
    ma rozwi\k{a}zanie $h(r) = \bigl(A\cos r + B\sin r\bigr)/r$.
    Cz\k{e}sto\'s\'c $\omega=1$, pó\l{}-okres $T_{\mathrm{half}}=\pi$.

  \item \textbf{Zanikanie amplitudy:} $|g(r)-1| \sim A_0/r$.
    Kolejne odbicia zachodz\k{a} przy $r_k \approx k\cdot\pi$,
    a amplituda przy $k$-tym odbiciu wynosi $A_0/(k\cdot\pi)$.

  \item \textbf{Warunek odbi\'c od ghost wall:}
    $A(r_k) \geq 1-g^*$, sk\k{a}d:
    \[
      \frac{g_0 - 1}{k\cdot\pi} \approx 1 - g^*
      \;\Longrightarrow\;
      \Delta g_0^* \;=\; \pi\,(1-g^*)
      \;=\; \pi\!\left(1 - e^{-1/(2\alpha)}\right).
    \]
\end{enumerate}
Dla $\alpha=2$:
\begin{equation}\label{eq:T4-deltag0-analytic}
  \boxed{
    \Delta g_0^*
    \;=\; \pi\,\bigl(1 - e^{-1/4}\bigr)
    \;\approx\; 0{,}695,
  }
\end{equation}
co odbiega od obserwowanego $0{,}675$ o~$3\%$
(reszta pochodzi z~$A_{\mathrm{tail}}/(g_0-1)\approx1{,}12$
oraz~$r_k \approx 1{,}7k\pi$ zamiast $k\pi$;
korekty cz\k{e}\'sciowo si\k{e} znosz\k{a}).
Weryfikacja numeryczna: \texttt{ex132} --- 14/14~PASS.
\end{proposition}

\begin{observation}[Zamkni\k{e}cie T-OP5]\label{obs:T4-top5-closed}
T-OP5 jest \textbf{zamkni\k{e}ty}.
Wzorzec liniowy~\eqref{eq:T4-bounce-thresholds} i~nachylenie
$\Delta g_0^* \approx 0{,}675$ s\k{a} konsekwencj\k{a} trzech elementów ODE:
(1)~cz\k{e}sto\'sci $\omega=1$ oscylacji liniowych przy $g=1$,
(2)~zaniku amplitudy ogona jak $1/r$,
(3)~g\l{}\'eboko\'sci ghost wall $1-g^* = 1-e^{-1/(2\alpha)}$.
Formu\l{}a $\Delta g_0^* = \pi(1-g^*)$ przewiduje nachylenie z~$3\%$
dok\l{}adno\'sci\k{a} bez parametr\'ow wolnych.
\end{observation}

% -------------------------------------------------------------------
\subsection{Podsumowanie --- status T-OP3}
\label{app:T4-summary}
% -------------------------------------------------------------------

\begin{center}
\small
\renewcommand{\arraystretch}{1.3}
\begin{tabular}{@{}p{5.5cm}lp{5cm}@{}}
\toprule
\textbf{Wynik} & \textbf{Status} & \textbf{\'Zr\'od\l{}o} \\
\midrule
BGH (monot.): $n_b$ ro\'snie wzdlu\.z drabiny
  & \statuslabel{Fakt num.}
  & \texttt{ex127}, B5, B9 \\[3pt]
Progi $g_0^*$: $1{,}650$, $2{,}250$, $2{,}900$, $3{,}549$, \ldots
  & \statuslabel{Fakt num.}
  & \texttt{ex127+ex130}, E1, E4 \\[3pt]
Wzorzec liniowy: $g_0^*(k) = 0{,}675k+1{,}582$
  & \statuslabel{Fakt num.}
  & \texttt{ex130}, E11 \\[3pt]
$g_0^{(\max)}(10\%) \approx g_0^*(3{\to}4)$ [$\Delta=0{,}147$]
  & \statuslabel{Twierdzenie}
  & \texttt{ex130}, E3 \\[3pt]
RCH: RMSE$/A(L_4) = 36{,}4\% > 15\%$
  & \statuslabel{Fakt num.}
  & \texttt{ex127}, B14 \\[3pt]
$g_0^{(\max)}(10\%) = 3{,}59 \in (g_0^\tau, g_0^{L_4})$
  & \statuslabel{Fakt num.}
  & \texttt{ex127}, B8 \\[3pt]
Tw.~\ref{thm:T4-rch}: TGP wyklucza $k=4$ wewn\k{e}trznie
  & \statuslabel{Twierdzenie}
  & \texttt{ex127} \\[3pt]
\eqref{eq:T4-chain}: RCH $\Rightarrow$ $N=3$ $\Rightarrow$ $Q_K=3/2$
  & \statuslabel{Twierdzenie}
  & ex127 + ex126 + ex118 \\[3pt]
$\Delta g_0^* = \pi(1-g^*)$ --- formu\l{}a analityczna
  & \statuslabel{Twierdzenie}
  & \texttt{ex132}, W1--W14 \\[3pt]
T-OP5: $g_0^{(\max)}$ i~$\Delta g_0^*$ --- \textbf{ZAMKNI\k{E}TE}
  & \statuslabel{Fakt num.+an.}
  & \texttt{ex130+ex132} \\[3pt]
RCH robustno\'s\'c (3~okna, 3~progi, Pearson)
  & \statuslabel{Fakt num.}
  & \texttt{ex128}, C1--C7 \\[3pt]
\l{}\k{a}\'ncuch RCH $\Rightarrow$ $N{=}3$ $\Rightarrow$ $Q_K{=}3/2$
  & \statuslabel{Twierdzenie}
  & \texttt{ex128}, C8--C16 \\
\bottomrule
\end{tabular}
\end{center}

\bigskip
\noindent\textbf{G\l{}\'owny wniosek.}
Mechanizm masy $m \propto A_{\mathrm{tail}}^4$ jest wiarygodny tylko
dla soliton\'ow z~$\mathrm{RMSE}/A < 10\%$, co ogranicza stabilne
genera\-cje do $k \leq 3$.  W~po\l{}\k{a}czeniu z~wynikami
dodatk\'ow T2 i~T3 (r\'ownanie~\ref{eq:T4-chain}),
\textbf{TGP PRZEWIDUJE $Q_K = 3/2$ z~pierwszych zasad}:
\begin{itemize}
  \item soliton nie mo\.ze stabilnie istnie\'c dla $k \geq 4$
    (kryterium RCH),
  \item wi\k{e}c $N = 3$ genera\-cje --- jedyna mo\.zliwo\'s\'c,
  \item wi\k{e}c $r = \sqrt{N-1} = \sqrt{2}$ (ex126),
  \item wi\k{e}c $Q_K = 2N/(N+1) = 3/2$.
\end{itemize}
Status: \statuslabel{Twierdzenie analityczne + weryfikacja numeryczna}
(\texttt{ex117--ex128}, 102/105~PASS \l{}\k{a}cznie).
Robustno\'s\'c RCH potwierdzona w~\texttt{ex128}
(16/16~PASS: 3~okna ogona, 3~progi RMSE, korelacja Pearsona
$[\text{ratio}=18{\times}]$, stosunek RMSE $[\text{ratio}=4{,}6{\times}]$).
